%%%%%%%%%%  scan idx 301  |  folio 288  %%%%%%%%%%
paraît surtout adaptée aux besoins de la Géométrie Algébrique et de la Géométrie Analytique, et une bonne définition du \(K_{\bullet}\) d'un espace topologique reste à trouver. Il semble d'ailleurs qu'il faudra sortir du cadre des complexes \(\mathcal O_X\)-linéaires, et définir pour un morphisme \(f:X\to Y\) (d'espaces topologiques, resp. de variétés mixtes (II App. II), etc.) une classe convenable de complexes de faisceaux vectoriels topologiques sur \(X\), devant englober par exemple, quand \(f\) est le morphisme structural d'une variété mixte, les complexes elliptiques sur \(X\) relativement à \(Y\). C'est ce que suggère le théorème de finitude (II App. II), ainsi qu'un théorème analogue (où \(Y\) est supposé réduit à un point) démontré par Grothendieck dans [3].

\textbf{2.12. Variances inhabituelles de \(K^{\bullet}\) et \(K_{\bullet}\).} Soit \(f:(X,A)\longrightarrow(Y,B)\) un morphisme de tor-dimension finie de topos annelés (III 3.1). Le foncteur \(Lf^*:D^-({}_B Y)\longrightarrow D^-({}_A X)\) envoie \(D^b({}_B Y)\) dans \(D^b({}_A X)\), donc (I 2.16.1) \(D^b({}_B Y)_{\coh}\) dans \(D^b({}_A X)_{\coh}\), et par suite définit un homomorphisme
\begin{equation}\tag{2.12.1}
 f^*:K_\bullet({}_B Y)\longrightarrow K_\bullet({}_A X).
\end{equation}
Si \(A\) et \(B\) sont commutatifs, la commutation de l'image inverse au produit tensoriel montre que l'homomorphisme (2.12.1) est compatible avec les structures de modules de \(K_{\bullet}(Y)\) et \(K_{\bullet}(X)\) sur \(K^{\bullet}(Y)\) et \(K^{\bullet}(X)\) respectivement (2.10), autrement dit que, pour \(a\in K^{\bullet}(Y)\) et \(y\in K_{\bullet}(Y)\), on a la formule
\begin{equation}\tag{2.12.2}
 f^*(ay)=a f^*(y).
\end{equation}
Si \(g:(Y,B)\to(Z,C)\) est un morphisme de tor-dimension finie de topos annelés, la transitivité des foncteurs images inverses implique que les homomorphismes (2.12.1) vérifient la formule :
\begin{equation}\tag{2.12.1.1}
 f^*g^*=(gf)^*.
\end{equation}

